\documentclass[11pt]{article}

\usepackage[margin=0.9in]{geometry}
\usepackage{graphicx}
\usepackage{amsmath}

\setlength{\parindent}{0pt}
\setlength{\parskip}{3pt}
\renewcommand{\baselinestretch}{0.95}

\title{\textbf{Deep Learning for Galaxy Spectral Super-Resolution}}
\author{Aryana Haghjoo}
\date{}

\begin{document}

\maketitle

\vspace{-0.4cm}

Low-resolution spectroscopy from JWST and the Nancy Grace Roman Space Telescope can observe large galaxy samples, but spectral features encoding key physics are blurred below instrumental resolution. I develop physics-informed deep learning frameworks that super-resolve low-resolution spectra by $\sim 10\times$, recovering hidden astrophysical information.

Using paired JWST/NIRSpec observations from the public JADES survey (same galaxies in high-resolution grating and low-resolution prism modes), I train generative models with physics-informed loss functions targeting emission line accuracy. Figure 1(a) demonstrates super-resolution on a $z=7.43$ galaxy, recovering spectral features invisible in prism data. Figure 1(b) shows the astrophysical impact: Lyman-alpha profiles depend sensitively on ionized bubble size ($R_b$). High-resolution profiles are distinguishable; prism-resolution profiles appear identical. Our method recovers this information, enabling physical measurements from low-resolution data.

\begin{figure}[h]
\centering
\includegraphics[width=0.95\textwidth]{figure1.pdf}
\vspace{-0.3cm}
\end{figure}

\vspace{-0.15cm}

\noindent
\textbf{Summer 2027 at IPAC:} In-person collaboration with Dr. Hemmati is essential. The fellowship will enable: (1) systematic validation of astrophysical diagnostic recovery using held-out JADES observations; (2) adaptation to Roman grism spectroscopy with IPAC's mission team; (3) development of multi-modal models integrating photometric and spectroscopic data. Direct access to IPAC's computational resources, mission specialists, and collaborators on related surveys accelerates model iteration and ensures our methodology leverages mission-specific characteristics. This immersive environment will significantly advance both research progress and quality compared to remote collaboration.

The work advances my dissertation, produces open-source software and trained models, and develops transferable methodology for future spectroscopic surveys.

\end{document}
